\chapter{Milestone M1.6: Fixed-Point Solvers, Grassmannian Geometry, and Multi-Basin Landscape}
\label{ch:m1_6}

% EDITORIAL QUESTION: How is the non-linear self-consistent mean-field problem formulated geometrically on the orthogonal Grassmannian, and what characterizes stationarity?
\section{Geometric Fixed-Point Formulation on the Orthogonal Grassmannian}
\label{sec:m1_6_grassmannian_geometry}

The self-consistent Hartree-Fock-Bogoliubov mean-field equations for the 4-flavor Majorana representation describe a nonlinear mapping on the orthogonal Grassmannian manifold bundle over the Brillouin zone:
\begin{equation}
\label{eq:m1_6_grassmannian_def}
\Gr_{\mathbb{C}}(2N_{\mathrm{sub}}, 4N_{\mathrm{sub}}) \equiv \left\{ \Pi \in \mathbb{C}^{4N_{\mathrm{sub}} \times 4N_{\mathrm{sub}}} \;\middle|\; \Pi^2 = \Pi = \Pi^\dagger, \; \operatorname{Tr}(\Pi) = 2N_{\mathrm{sub}} \right\},
\end{equation}
subject to the particle-hole constraint $\Pi^-(-\mathbf{k}) = \Id - (\Pi^-(\mathbf{k}))^*$. For the 2-sublattice primitive cell ($N_{\mathrm{sub}} = 2$), the fiber dimension is $4N_{\mathrm{sub}} = 8$ and each occupied projector has trace $\operatorname{Tr}(\Pi^-) = 4$.
At zero temperature ($T=0$), the physical two-point correlator $P(\mathbf{k}) = 2\Pi^-(\mathbf{k})$ projects onto the negative-energy eigenspace of the momentum-space Hamiltonian $\mathcal{H}(\mathbf{k}) \equiv \frac{\ii}{4} M(\mathbf{k})$.

\begin{theorem}[Grassmannian Stationarity and KKT Commutator Condition]
\label{thm:m1_6_kkt_condition}
A mean-field state $X^* = (\mathrm{U}^*, \mathrm{V}^*, \boldsymbol{\lambda}^*)$ is a self-consistent Karush-Kuhn-Tucker (KKT) stationary point if and only if two conditions hold simultaneously:
\begin{enumerate}
\item \textbf{Spectral Commutator Stationarity:} At every wavevector $\mathbf{k} \in \mathrm{BZ}$, the single-particle density matrix $\Pi^-(\mathbf{k})$ commutes with the assembled Hamiltonian:
\begin{equation}
\label{eq:m1_6_kkt_comm}
\epsilon_{\mathrm{comm}}(\mathbf{k}) \equiv \| [\mathcal{H}(\mathbf{k}; X^*), \Pi^-(\mathbf{k})] \|_F = 0, \qquad \epsilon_{\mathrm{comm}} \equiv \frac{1}{N_{\mathbf{k}}} \sum_{\mathbf{k} \in \mathrm{BZ}} \epsilon_{\mathrm{comm}}(\mathbf{k}) = 0.
\end{equation}
\item \textbf{Gauge Constraint Admissibility:} On every physical site $s \in \{0, \dots, N_{\mathrm{sub}}-1\}$ and for all three gauge flavors $\gamma \in \{1, 2, 3\}$, the gauge residual vanishes:
\begin{equation}
\label{eq:m1_6_gauge_condition}
g_s^\gamma \equiv \operatorname{Tr}(\mathrm{G}^\gamma \mathrm{V}_s^*) = 0 \implies \epsilon_{\mathrm{gauge}} \equiv \|\mathbf{g}\|_\infty = 0.
\end{equation}
\end{enumerate}
\end{theorem}
\begin{proof}
The variational Hartree-Fock ground-state energy functional is $E_{\mathrm{MF}}[P, \boldsymbol{\lambda}] = \langle \mathcal{H}[P] \rangle_P - E_{\mathrm{exchange}}[P] - \sum_{s,\gamma} \lambda_s^\gamma g_s^\gamma$. The tangent space to the Grassmannian at $\Pi^-$ is parameterized by skew-Hermitian commutators $\delta \Pi = [\Delta, \Pi^-]$ where $\Delta^\dagger = -\Delta$. First-order stationarity $\delta E_{\mathrm{MF}} = 0$ requires $\operatorname{Tr}(\mathcal{H}(\mathbf{k}) [\Delta(\mathbf{k}), \Pi^-(\mathbf{k})]) = 0$ for all skew-Hermitian variations $\Delta(\mathbf{k})$, which holds if and only if $[\mathcal{H}(\mathbf{k}), \Pi^-(\mathbf{k})] = \mathbf{0}$. Stationarity with respect to the Lagrange multipliers $\lambda_s^\gamma$ yields the gauge condition $\partial E / \partial \lambda_s^\gamma = -g_s^\gamma / 4 = 0$.
\end{proof}

% EDITORIAL QUESTION: What are the mathematical foundations of fixed-point iterations, and how do contractivity and asymptotic convergence rates govern iterative solver design?
\section{Foundations of Fixed-Point Iteration and Optimization Theory}
\label{sec:m1_6_foundations_fixed_point}

Before analyzing accelerated numerical solvers, we establish the general mathematical framework for solving non-linear algebraic systems without relying on algorithm-specific jargon.

\subsection{Banach Contraction Principle and Asymptotic Rates}
Let $\mathcal{X}$ denote a real Banach space equipped with norm $\|\cdot\|$ and dimension $D$. A self-consistent mean-field mapping $\Phi: \mathcal{X} \to \mathcal{X}$ defines a root-finding problem for the residual map:
\begin{equation}
G(x) \equiv \Phi(x) - x = 0.
\end{equation}
A point $x^* \in \mathcal{X}$ satisfying $G(x^*) = 0 \iff \Phi(x^*) = x^*$ is a fixed point.

\begin{definition}[Contraction Factor and Convergence Classes]
Let $\{x_k\}_{k=0}^\infty$ be a sequence converging to $x^*$, and let $e_k \equiv x_k - x^*$ denote the error vector at iteration $k$. The asymptotic convergence rate is classified by the limit:
\begin{equation}
\lim_{k \to \infty} \frac{\|e_{k+1}\|}{\|e_k\|^p} = C.
\end{equation}
\begin{enumerate}
\item \textbf{Linear Convergence ($p=1, C \in (0, 1)$):} The error decreases geometrically, $\|e_{k+1}\| \le \sigma \|e_k\|$, with asymptotic contraction factor $\sigma = \rho(\nabla \Phi(x^*)) < 1$, where $\rho(\cdot)$ is the spectral radius.
\item \textbf{Superlinear Convergence ($p=1, C=0$, or $1 < p < 2$):} The step-to-step contraction ratio $\sigma_k \equiv \|e_{k+1}\| / \|e_k\| \to 0$ as $k \to \infty$.
\item \textbf{Quadratic Convergence ($p=2, C > 0$):} The number of correct significant digits doubles at every single iteration step: $\|e_{k+1}\| \le C \|e_k\|^2$.
\end{enumerate}
\end{definition}

\subsection{Plain Fixed-Point and Damped Picard Iteration}
The simplest strategy is the unaccelerated fixed-point iteration:
\begin{equation}
x_{k+1} = \Phi(x_k).
\end{equation}
By the Banach Fixed-Point Theorem, this sequence converges locally if and only if $\Phi$ is contractive in a neighborhood of $x^*$, which requires the spectral radius of the Jacobian to satisfy $\rho(\nabla \Phi(x^*)) < 1$. When $\rho(\nabla \Phi) \ge 1$, the iteration diverges geometrically.

To stabilize unstable fixed points or expand the basin of attraction, \textbf{Damped Picard iteration} mixes the previous iterate with the evaluated map using a convex combination parameter $\alpha \in (0, 1]$:
\begin{equation}
x_{k+1} = (1 - \alpha) x_k + \alpha \Phi(x_k) = x_k + \alpha G(x_k).
\end{equation}
The effective iteration matrix is $T_\alpha = (1 - \alpha) I + \alpha \nabla \Phi(x^*)$. If all eigenvalues of the residual Jacobian $J_G = \nabla \Phi - I$ have negative real parts, there exists a sufficiently small damping $\alpha > 0$ such that $\rho(T_\alpha) < 1$, guaranteeing linear convergence. However, as the system size increases or near critical phase boundaries where $\rho \to 1^-$, damped Picard suffers from extreme \textit{critical slowing down}, requiring hundreds of iterations.

\subsection{Newton-Raphson Method and the Jacobian Curse}
To achieve quadratic convergence ($p=2$), the classical Newton-Raphson method expands the residual $G(x)$ to first order in a Taylor series:
\begin{equation}
G(x_k + \delta x) \approx G(x_k) + J_G(x_k) \delta x = 0 \implies \delta x_k = - [J_G(x_k)]^{-1} G(x_k),
\end{equation}
yielding the update $x_{k+1} = x_k + \delta x_k$.

While Newton's method converges quadratically near $x^*$, it suffers from the \textbf{Jacobian Curse} in quantum many-body physics:
\begin{enumerate}
\item \textbf{Assembly Cost:} In our 4-flavor Majorana problem, assembling the exact $D \times D$ Jacobian $J_G$ requires $D$ directional evaluations of the full Brillouin-zone eigensolver, costing $\mathcal{O}(D \cdot M^3 N_{\mathbf{k}})$ FLOPs per step.
\item \textbf{Inversion Cost:} Solving the linear system $J_G \delta x = -G$ requires dense LU factorization costing $\frac{2}{3} D^3$ FLOPs. For large supercells ($D > 10^4$), storing and factoring $J_G$ is computationally impossible.
\end{enumerate}

% EDITORIAL QUESTION: How do Quasi-Newton methods (Good Broyden, Bad Broyden, DIIS, L-Broyden) circumvent Jacobian assembly while maintaining superlinear contraction?
\section{Quasi-Newton Methods: Broyden's Zoo, DIIS, and JFNK}
\label{sec:m1_6_quasi_newton_zoo}

Quasi-Newton methods replace the true Jacobian $J_G(x_k)$ (or its inverse $H_k \approx [J_G(x_k)]^{-1}$) with an approximation that is updated at each step using only the trajectory vectors:
\begin{equation}
s_k \equiv x_{k+1} - x_k, \qquad y_k \equiv G(x_{k+1}) - G(x_k).
\end{equation}
Every valid quasi-Newton update must satisfy the \textbf{Secant Equation}:
\begin{equation}
B_{k+1} s_k = y_k \iff H_{k+1} y_k = s_k.
\end{equation}

\subsection{Good Broyden (Type-I) vs. Bad Broyden (Type-II)}
Because the secant equation under-determines the $D \times D$ matrix $B_{k+1}$ (1 vector constraint on $D^2$ entries), an optimization criterion is required:
\begin{enumerate}
\item \textbf{Good Broyden (Type-I):} Minimizes the change in the direct Jacobian: $\min \|B_{k+1} - B_k\|_F$ subject to $B_{k+1} s_k = y_k$. Applying the Sherman-Morrison inversion lemma yields the explicit rank-1 update on the inverse Jacobian $H_k$:
\begin{equation}
\label{eq:m1_6_good_broyden_def}
H_{k+1} = H_k + \frac{s_k - H_k y_k}{s_k^T H_k y_k} s_k^T H_k.
\end{equation}
It is called "Good" Broyden because it preserves orthogonality properties in linear spaces and avoids direct matrix inversion.
\item \textbf{Bad Broyden (Type-II):} Minimizes the change directly in the inverse Jacobian: $\min \|H_{k+1} - H_k\|_F$ subject to $H_{k+1} y_k = s_k$, which produces:
\begin{equation}
\label{eq:m1_6_bad_broyden_def}
H_{k+1} = H_k + \frac{s_k - H_k y_k}{y_k^T y_k} y_k^T.
\end{equation}
While simpler, Type-II updates are known in numerical analysis to be less robust for general non-linear systems than Type-I updates, though they form the conceptual basis for multi-secant methods.
\end{enumerate}

\subsection{Joint DIIS / Anderson Acceleration (Multi-Secant)}
Direct Inversion in the Iterative Subspace (DIIS), mathematically equivalent to Anderson Acceleration, expands the trial state over the history of the $m$ most recent iterates $\{x_{k-j}\}_{j=0}^{m-1}$:
\begin{equation}
x_{k+1} = \sum_{j=0}^{m-1} c_j \Phi(x_{k-j}) = \sum_{j=0}^{m-1} c_j (x_{k-j} + G_{k-j}),
\end{equation}
where the coefficients $\mathbf{c} = (c_0, \dots, c_{m-1})^T$ minimize the Euclidean norm of the composite residual $\| \sum c_j G_{k-j} \|_2^2$ subject to the partition-of-unity constraint $\sum_{j=0}^{m-1} c_j = 1$.
This yields the $(m+1) \times (m+1)$ Karush-Kuhn-Tucker (KKT) linear system:
\begin{equation}
\begin{pmatrix}
B & -\mathbf{1} \\
-\mathbf{1}^T & 0
\end{pmatrix}
\begin{pmatrix}
\mathbf{c} \\
\mu
\end{pmatrix}
=
\begin{pmatrix}
\mathbf{0} \\
-1
\end{pmatrix},
\qquad
B_{pq} = \langle G_{k-p}, G_{k-q} \rangle + \epsilon_{\mathrm{reg}} \delta_{pq} \max(\|G_{k-p}\|^2, 10^{-14}).
\end{equation}
Because $m \ll D$ (typically $m \in [6, 10]$), factoring the small $m \times m$ Gram matrix $B$ takes microseconds, making DIIS an extremely powerful multi-secant accelerator.

\subsection{Limited-Memory Broyden (L-Broyden)}
In high dimensions ($D \gg 10^3$), storing the $D \times D$ matrix $H_k$ requires $\mathcal{O}(D^2)$ storage and $\mathcal{O}(D^2)$ work.
L-Broyden stores only the $m$ most recent secant pairs $\{(s_j, y_j)\}_{j=0}^{m-1}$, evaluating the quasi-Newton step $H_k G_k$ via recursive vector inner products matching the Sherman-Morrison sequence.
The memory footprint is reduced to strictly $\mathcal{O}(mD)$ and the step overhead to $\mathcal{O}(m^2 D)$, with zero matrix storage.

\subsection{Jacobian-Free Newton-Krylov (JFNK)}
Instead of building an approximate Jacobian, JFNK solves the exact Newton step $J_G(x_k) \delta x = -G(x_k)$ using the GMRES (Generalized Minimal Residual) Krylov subspace iteration. GMRES requires only matrix-vector products $J_G v$, which are evaluated via finite differences:
\begin{equation}
J_G(x) v \approx \frac{G(x + \epsilon v) - G(x)}{\epsilon}, \qquad \epsilon = \epsilon_0 \frac{1 + \|x\|_2}{\|v\|_2}.
\end{equation}
Each GMRES inner step requires one additional evaluation of the physical mean-field map $\Phi(x + \epsilon v)$. JFNK achieves true quadratic-to-superlinear convergence with zero Jacobian storage, at the cost of multiple map evaluations per outer step.

\subsection{Concrete Leading-Order Computational Complexity}
The total computational cost per iteration of each solver paradigm is summarized with exact leading-order operational prefactors in Table~\ref{tab:m1_6_solver_leading_order_flops}.

\begin{table}[t]
\centering
\small
\caption{Exact leading-order arithmetic FLOP formulas and memory footprints per iteration step.}
\label{tab:m1_6_solver_leading_order_flops}
\begin{tabular}{lcccc}
\toprule
\textbf{Solver Paradigm} & \textbf{Map Evals / Step} & \textbf{Solver Overhead FLOPs} & \textbf{Dominant Evaluation FLOPs ($\Phi$)} & \textbf{Memory Footprint} \\
\midrule
\textbf{Plain FP / Picard} & $1$ & $2D$ & $\frac{16}{3} M^3 N_{\mathbf{k}}$ & $2 D \times 8\text{ B}$ \\
\textbf{Joint DIIS ($m$)} & $1$ & $2mD + \frac{1}{3}m^3$ & $\frac{16}{3} M^3 N_{\mathbf{k}}$ & $2m D \times 8\text{ B}$ \\
\textbf{Good Broyden (Full)} & $1$ & $2D^2 + 4D$ & $\frac{16}{3} M^3 N_{\mathbf{k}}$ & $D^2 \times 8\text{ B}$ \\
\textbf{L-Broyden ($m$)} & $1$ & $4m^2 D + 2mD$ & $\frac{16}{3} M^3 N_{\mathbf{k}}$ & $2m D \times 8\text{ B}$ \\
\textbf{JFNK ($k_{\mathrm{gmres}}$)} & $1 + k_{\mathrm{gmres}}$ & $2 k_{\mathrm{gmres}} D$ & $(1 + k_{\mathrm{gmres}}) \cdot \frac{16}{3} M^3 N_{\mathbf{k}}$ & $k_{\mathrm{gmres}} D \times 8\text{ B}$ \\
\midrule
\textbf{Real Skew Eigensolve} & $1$ & -- & $\left(\frac{4}{3} M^3 + 6 M^3\right) N_{\mathbf{k}}$ & $\frac{1}{2} M^2 \times 8\text{ B}$ \\
\textbf{Matrix Sign Purif ($k_{\mathrm{NS}}$)} & $1$ & -- & $4 k_{\mathrm{NS}} M^3$ (Real \texttt{dgemm}) & $M^2 \times 8\text{ B}$ \\
\bottomrule
\end{tabular}
\end{table}

\begin{proposition}[Convergence and Contraction Telemetry]
\label{prop:m1_6_convergence_rates}
The fixed-point contraction factor at step $k$ is defined by:
\begin{equation}
\label{eq:m1_6_contraction_factor}
\mu_k \equiv \frac{\|R_{k+1}\|_F}{\|R_k\|_F}.
\end{equation}
Under unperturbed seeds, Joint DIIS ($m=6$) and Good Broyden achieve superlinear-to-geometric convergence ($\mu_k \le 0.01$) requiring only 8--13 map evaluations to reach $\epsilon_{\mathrm{comm}} < 10^{-10}$. In contrast, plain Fixed Point and Picard exhibit sluggish linear convergence ($\mu_k \approx 0.99$), stagnating at $\sim 10^{-6}$ after 60 iterations due to the disparity between rapid fermionic density relaxation and slow fixed-step dual gradient ascent ($\eta_\lambda = 0.2$).
\end{proposition}

% EDITORIAL QUESTION: Why is skew-symmetric retraction necessary during fixed-point updates, and how was it verified through A/B testing?
\section{Skew-Symmetric Retraction and A/B Testing}
\label{sec:m1_6_ab_retraction}

The physical real-space covariance matrices $V_s \in \mathbb{R}^{4 \times 4}$ and bond matrices $U_b \in \mathbb{R}^{4 \times 4}$ belong strictly to the Lie algebra $\mathfrak{so}(4)$ and the affine manifold of CAR correlators. During numerical iterative updates, small floating-point roundoff errors can leak into the symmetric sector.

To systematically test the necessity of skew-symmetric projection, a 4-arm A/B benchmark was executed in \texttt{doc/rnd/phase\_1/m1\_6\_t2\_skew\_symmetry\_ab\_benchmark.md}:
\begin{itemize}
\item \textbf{Arm A (Raw Unprojected):} Leaves iterates unconstrained.
\item \textbf{Arm B (Output Symmetrization):} Enforces $V_s \leftarrow \frac{1}{2}(V_s - V_s^T)$ only when evaluating the Hamiltonian.
\item \textbf{Arm C (Post-Step Retraction):} Projects $X \leftarrow \mathcal{P}_{\mathfrak{so}}(X)$ immediately after every accelerator step.
\item \textbf{Arm D (Retraction + Symmetrization):} Dual enforcement.
\end{itemize}

\begin{proposition}[Failure of Unprojected Iteration under Krylov Subspaces]
\label{thm:m1_6_ab_failure}
In Arm A, finite-difference perturbations inside GMRES Arnoldi loops accumulate asymmetric roundoff error exceeding domain validation thresholds ($\|V_s + V_s^T\|_\infty > 10^{-12}$), causing hard solver crashes. In contrast, Arms B and C produce numerically identical trajectories to machine precision ($< 10^{-14}$) while enforcing $\|V_s + V_s^T\|_\infty = 0$ with 0\% computational overhead.
\end{proposition}

% EDITORIAL QUESTION: What is the global structure of the multi-basin landscape under field h || [111], and where does the thermodynamic transition occur?
\section{The Multi-Basin Landscape and Hysteresis under Field $\mathbf{h} \parallel [111]$}
\label{sec:m1_6_hysteresis}

Under an external Zeeman field $\mathbf{h} = \frac{h}{\sqrt{3}}(1, 1, 1)$, the self-consistent mean-field energy functional $E_{\mathrm{MF}}[\Omega]$ exhibits a multi-basin landscape with competing local minima:
\begin{enumerate}
\item \textbf{Topological Quantum Spin Liquid (QSL) Basin:} Characterized by extensive flux $W_p \in [0.53, 1.0]$, vanishing magnetization at $h=0$, and preserving exact $C_3$ rotational symmetry.
\item \textbf{Field-Polarized Basin:} Characterized by zero plaquette flux ($W_p = 0.0000$), saturated magnetization $m_x = m_y = m_z = 1/\sqrt{3}$, and exact classical product energy:
\begin{equation}
\label{eq:m1_6_classical_energy}
E_{\mathrm{classical}}(h) = -h - \frac{1}{2}.
\end{equation}
\end{enumerate}

\begin{proposition}[First-Order Thermodynamic Transition and Metastability Window]
\label{thm:m1_6_critical_transition}
A high-resolution zoom sweep ($\Delta h = 0.002$) across $h \in [0.280, 0.520]$ establishes:
\begin{enumerate}
\item \textbf{Thermodynamic Crossing:} The absolute global minimum switches at:
\begin{equation}
\label{eq:m1_6_hc_star}
h_c^* = 0.3325 \pm 0.0002.
\end{equation}
For $h < h_c^*$, $E_{\mathrm{QSL}}(h) < E_{\mathrm{pol}}(h)$, establishing the QSL as the true physical ground state.
\item \textbf{Forward Spinodal Breakdown:} Under forward adiabatic continuation from $h=0$, the QSL branch remains metastable for $h \in [h_c^*, h_{c2}^+]$ until the forward spinodal:
\begin{equation}
\label{eq:m1_6_hc2_plus}
h_{c2}^+ = 0.400 \pm 0.002.
\end{equation}
At $h = 0.400$, the fixed-point Jacobian loses contractivity ($\mu_k \to 1^-$), leading to critical slowing down and an irreversible jump into the polarized fixed point.
\item \textbf{Backward Spinodal Breakdown:} Under backward continuation from $h=5.0$, the polarized state remains a local minimum all the way down to $h_{c1}^- = 0.0$, generating a broad hysteresis loop.
\end{enumerate}
\end{proposition}

\begin{figure}[htbp]
\centering
\begin{tikzpicture}
\begin{axis}[
    labbook standard,
    xlabel={Magnetic Field $h/|K|$ ($\mathbf{h} \parallel [111]$)},
    ylabel={Energy per Site $E_{\mathrm{site}}$},
    legend pos=south west,
    xmin=0.28, xmax=0.52,
    ymin=-1.02, ymax=-0.81
]
    \addplot[labbookteal, thick]
        table[x index=0, y index=1] {../../data/processed/benchmarks/m1_6/critical_zoom_plot.dat};
    \addlegendentry{QSL Forward Branch ($W_p > 0$)}

    \addplot[labbookcrimson, dashed, thick]
        table[x index=0, y index=2] {../../data/processed/benchmarks/m1_6/critical_zoom_plot.dat};
    \addlegendentry{Polarized Backward Branch ($W_p = 0$)}

    \draw[dotted, thick, black] (axis cs:0.3325,-1.02) -- (axis cs:0.3325,-0.81)
        node[pos=0.85, right, font=\footnotesize] {$h_c^* = 0.3325$ (Crossing)};

    \draw[dashed, thick, labbookamber] (axis cs:0.4000,-1.02) -- (axis cs:0.4000,-0.81)
        node[pos=0.15, right, font=\footnotesize] {$h_{c2}^+ = 0.400$ (Spinodal)};
\end{axis}
\end{tikzpicture}
\caption{High-resolution critical region zoom ($\Delta h = 0.002$) showing the first-order thermodynamic crossing $h_c^* \approx 0.3325$ and forward spinodal breakdown $h_{c2}^+ \approx 0.400$ on the orthogonal Grassmannian manifold $\Gr(4, 8)$.}
\label{fig:m1_6_critical_zoom}
\end{figure}


% EDITORIAL QUESTION: How does the topological Majorana gap open under a [111] field, and does mean-field theory recover Kitaev's chiral perturbation scaling?
\section{Topological Majorana Gap Opening and Chiral Perturbation Scaling}
\label{sec:m1_6_majorana_gap_scaling}

In the unperturbed Kitaev honeycomb model ($h = 0$), the zero-flux sector features gapless Majorana Dirac cones at the high-symmetry points $K = (\frac{1}{3}, \frac{2}{3})$ and $K' = (\frac{2}{3}, \frac{1}{3})$. In perturbation theory, Kitaev proved that an external field $\mathbf{h} \parallel [111]$ generates an effective time-reversal-breaking three-spin interaction at third order:
\begin{equation}
\label{eq:m1_6_kitaev_kappa}
H^{(3)} = -\kappa \sum_{\langle\langle j, k\rangle\rangle} \sigma_j^\alpha \sigma_l^\gamma \sigma_k^\beta, \qquad \kappa \approx \frac{h_x h_y h_z}{\Delta_{\mathrm{flux}}^2} \propto \frac{h^3}{|K|^2},
\end{equation}
which opens a topological Dirac mass gap $\Delta_K \approx 6\sqrt{3}\kappa \propto h^3$.

\begin{theorem}[Algebraic Exponent Selection Rule for Dirac Mass Opening]
\label{thm:m1_6_cubic_gap}
Let $\Delta_K(\mathbf{h}) \equiv \varepsilon_5(K) - \varepsilon_4(K)$ denote the single-particle excitation gap at the Dirac point in self-consistent 4-flavor Majorana mean-field theory under spatial inversion, threefold lattice rotation $C_3$, and time-reversal $\mathcal{T}$:
\begin{enumerate}
\item Linear $\mathcal{O}(h)$ and quadratic $\mathcal{O}(h^2)$ scalar Dirac mass contributions vanish identically due to the combination of particle-hole symmetry $\mathcal{H}(-\mathbf{k}) = -\mathcal{H}^*(\mathbf{k})$, time-reversal parity, and discrete $C_3$ rotational invariance.
\item The lowest-order non-vanishing $\mathcal{T}$-odd, $C_3$-invariant scalar polynomial in the field components $(h_x, h_y, h_z)$ is the chiral triple product $\mathcal{I}^{(3)}(\mathbf{h}) = h_x h_y h_z$. Along $\mathbf{h} = \frac{h}{\sqrt{3}}(1, 1, 1)$, $\mathcal{I}^{(3)}(\mathbf{h}) = \frac{h^3}{3\sqrt{3}}$, algebraically mandating a cubic asymptotic gap opening $\Delta_K(h) = c_{\mathrm{MF}} (h/|K|)^3 + \mathcal{O}(h^5)$.
\end{enumerate}
\end{theorem}
\begin{proof}
By time-reversal invariance $\mathcal{T}$, the single-particle Bogoliubov-de Gennes generator satisfies $\mathcal{T}\mathcal{H}(\mathbf{k}, \mathbf{h})\mathcal{T}^{-1} = \mathcal{H}(-\mathbf{k}, -\mathbf{h})$, mapping $K \leftrightarrow -K \equiv K'$. Any scalar Dirac mass term $m_D \tau_z$ opening a gap at $K$ with opposite sign at $K'$ must be odd under $\mathcal{T}$ ($\mathbf{h} \mapsto -\mathbf{h}$) and invariant under the threefold point-group rotation $C_3$ that permutes field components and bond axes cyclically $(x \to y \to z \to x)$. Quadratic polynomials $h_\alpha h_\beta$ are $\mathcal{T}$-even and can only shift the Dirac node position in momentum space without opening a mass gap. Linear combinations $\sum_\alpha a_\alpha h_\alpha$ transform as a vector under $C_3$ except for the symmetric sum $h_x + h_y + h_z$, which acts identically on both sublattices A and B and therefore preserves sublattice chiral parity at $K$. Consequently, the lowest-order polynomial invariant that is simultaneously $\mathcal{T}$-odd and $C_3$-symmetric is:
\begin{equation}
\mathcal{I}^{(3)}(\mathbf{h}) = h_x h_y h_z = \frac{h^3}{3\sqrt{3}},
\end{equation}
establishing that the leading Dirac mass gap opens strictly at cubic order $\mathcal{O}(h^3)$.
\end{proof}

\begin{proposition}[Empirical Cubic Dirac Mass Scaling in 4-Flavor Mean-Field Theory]
\label{prop:m1_6_cubic_gap_fit}
High-resolution numerical evaluation across $h \in [0.005, 0.150]$ on a $24 \times 24$ Brillouin-zone mesh using \texttt{scripts/m1\_6/benchmark\_majorana\_gap.py} establishes:
\begin{enumerate}
\item The single-particle Bloch gap $\Delta_K(h)$ follows the cubic power-law regression:
\begin{equation}
\label{eq:m1_6_cubic_fit}
\Delta_K(h) = (2.5921 \pm 0.0034) \cdot \left(\frac{h}{|K|}\right)^{2.9936 \pm 0.0006} \approx 2.5921 \left(\frac{h}{|K|}\right)^3,
\end{equation}
with coefficient of determination $R^2 > 0.99999$.
\item The Brillouin-zone corner $K$ remains the true global minimum bandgap throughout the weak-field regime: $\Delta_K(h) = \min_{\mathbf{k}} \Delta(\mathbf{k})$.
\end{enumerate}
\end{proposition}

\begin{figure}[htbp]
\centering
\begin{tikzpicture}
\begin{loglogaxis}[
    labbook standard,
    xlabel={Magnetic Field $h/|K|$ ($\mathbf{h} \parallel [111]$)},
    ylabel={Majorana Dirac Gap $\Delta_K(h)$},
    legend pos=north west,
    xmin=0.004, xmax=0.4,
    ymin=1e-7, ymax=0.2
]
    \addplot[labbookteal, thick, mark=*, mark size=2pt]
        table[x index=0, y index=1] {../../data/processed/benchmarks/m1_6/majorana_gap_plot.dat};
    \addlegendentry{Measured Dirac Gap $\Delta_K(h)$}

    \addplot[labbookcrimson, dashed, thick]
        table[x index=0, y index=3] {../../data/processed/benchmarks/m1_6/majorana_gap_plot.dat};
    \addlegendentry{Cubic Reference: $2.592 \cdot (h/|K|)^3$}
\end{loglogaxis}
\end{tikzpicture}
\caption{Topological Majorana mass gap opening at the Dirac point $K$. The single-particle gap follows the strictly cubic power law $\Delta_K(h) \approx 2.592 h^3$, reflecting the cancellation of linear and quadratic on-site tilts and recovering Kitaev's third-order chiral perturbation $\kappa \propto h^3/|K|^2$.}
\label{fig:m1_6_majorana_gap}
\end{figure}


\subsection{Reconciling Perturbative Vison Scaling with Self-Consistent Mean-Field Theory (K-019)}
\label{sec:m1_6_perturbative_reconciliation}
% CLAIM: K-019 (provenance: doc/rnd/phase_1/m1_10_t1_exact_perturbation_and_normalization_reconciliation.md, test: test/unit/test_post_spinodal_kkt_audit.py::test_hc_p1_01_perturbative_gap_reconciliation_identities)

A central theoretical question (\texttt{HC-P1-01}) is how the self-consistent mean-field prefactor $c_{\mathrm{MF}} = 2.5921 \pm 0.0034$ relates to third-order spin perturbation theory on the physical gauge-singlet subspace. Operator-level structural analysis separates this comparison into three distinct components:
\begin{enumerate}
\item \textbf{Canonical Bogoliubov Factor of 4 (Exact Normalization Identity):} From Proposition~\ref{prop:m1_2_spectral_symmetries}, the physical fermionic quasiparticle excitation spectrum (eigenvalues of $\ii A$) satisfies $\varepsilon_n(\mathbf{k}) = 4\mathcal{E}_n(\mathbf{k})$ relative to the code's Bloch matrix $H_{\mathrm{Bloch}}(\mathbf{k}) = \frac{\ii}{4}A(\mathbf{k})$. In canonical Majorana normalization, the self-consistent mean-field gap is therefore:
\begin{equation}
\label{eq:m1_6_canonical_mf_gap}
\Delta_{\mathrm{canonical}}^{\mathrm{MF}}(h) = 4 \Delta_K^{\mathrm{MF}}(h) = 4 \times 2.5921 \left(\frac{h}{|K|}\right)^3 = 10.3684 \frac{h^3}{|K|^2}.
\end{equation}
\item \textbf{Feshbach Two-Vison Continuum vs. Local Gauge-Band Energy Scales:} In exact spin perturbation theory on the physical subspace $\mathcal{F}_{\mathrm{phys}}$, acting with a single Pauli operator $\sigma_i^\gamma = \ii c_i^0 c_i^\gamma$ flips the adjacent $\mathbb{Z}_2$ bond gauge link $u_{\langle ij\rangle_\gamma} \to -u_{\langle ij\rangle_\gamma}$ and nucleates a pair of localized $\mathbb{Z}_2$ visons together with dispersive matter Majorana excitations. In Kitaev's static single-mode approximation (replacing the full two-vison + matter Feshbach resolvent $Q(E_0 - H_K)^{-1}Q$ by the isolated two-vison threshold $\Delta_{\mathrm{flux}} \approx 0.264 |K|$), the third-order chiral hopping amplitude $\kappa_{\mathrm{static}} \approx h_x h_y h_z / \Delta_{\mathrm{flux}}^2$ yields a static reference scale $\Delta_{\mathrm{static}} = 12\sqrt{3}\kappa_{\mathrm{static}} = \frac{4}{\Delta_{\mathrm{flux}}^2} h^3 \approx 57.40 h^3/|K|^2$ (the exact continuum-integrated Feshbach value $\kappa_{\mathrm{exact}}$ is derived independently in the exact perturbation track). By contrast, in unprojected Gaussian mean-field theory, background $\mathbb{Z}_2$ fluxes are not flipped locally; instead, the intermediate single-particle excitations involve the localized flat gauge-Majorana bands with two-particle gap $2\Delta_{\mathrm{gauge}} = 2|K|\bar{U} \approx 1.04973 |K|$ (\texttt{K-004}, \texttt{K-006}). The bare squared-denominator ratio between these two intermediate scales is:
\begin{equation}
\label{eq:m1_6_denominator_ratio}
\left(\frac{2\Delta_{\mathrm{gauge}}}{\Delta_{\mathrm{flux}}}\right)^2 = \left(\frac{1.0497315}{0.2640}\right)^2 = 15.8106.
\end{equation}
\item \textbf{Structural Role of Self-Consistent KKT Multipliers and Fock Dressing (Contract H6):} Unlike spin perturbation theory---where three successive Zeeman applications $\sigma_j^\alpha \sigma_k^\gamma \sigma_l^\beta$ directly contract the internal gauge fermions via the local projection identity $D_k = -c_k^0 c_k^x c_k^y c_k^z = +1$---in 4-flavor HFB mean-field theory the bare Zeeman matrix $-\sum_\gamma h_\gamma N^\gamma$ only couples the dispersive matter flavor $c^0$ to the gauge flavors $c^\gamma$ ($N^\gamma_{0\gamma} = -1$) and contains no direct $c^\alpha \leftrightarrow c^\beta$ mixing. Consequently, a naive substitution of $2\Delta_{\mathrm{gauge}}$ into the spin-perturbation formula $\frac{4}{\Delta E^2}h^3$ is structurally invalid. Instead, the $\mathcal{O}(h^3)$ chiral Dirac mass in HFB theory is generated by the \textbf{self-consistent linear response} at $\mathcal{O}(h)$: the field induces both the KKT Lagrange multiplier counter-field $\lambda_s^\gamma = \mathcal{O}(h)$---which multiplies the dual gauge generators $G^\gamma$ ($G^\gamma_{\alpha\beta} = -\epsilon_{\gamma\alpha\beta}$) that directly mix distinct gauge flavors $c^\alpha \leftrightarrow c^\beta$---and the off-diagonal Fock bond self-energy $\delta \mathrm{U}_b = \mathcal{O}(h)$. Expanding the KKT stationarity condition $[\mathcal{H}(\mathbf{k}), \Pi^-(\mathbf{k})] = 0$ analytically to third order in $h$ to obtain the closed-form prefactor $c_{\mathrm{MF}} = 2.59214\dots$ is formalized in Contract~\ref{con:h6_analytical_cubic_prefactor}.
\end{enumerate}

\subsection{Post-Spinodal KKT Saddle Excursions and Coexistence Domain Walls (K-023)}
\label{sec:m1_6_post_spinodal_kkt_audit}
% CLAIM: K-023 (provenance: scripts/m1_9/audit_post_spinodal_and_master_table.py, test: test/unit/test_post_spinodal_kkt_audit.py::test_hc_p1_02_post_spinodal_kkt_admissibility_at_h0450)

In the raw 60-iteration forward sweep of Figure~\ref{fig:m1_6_critical_zoom}, the QSL continuation trace exhibits serrated downward energy dips across $h \in [0.400, 0.454]$ (reaching $E_{\mathrm{site}} = -1.1904$ and $m_{111} = 1.1249 > 1$ at $h = 0.450$). A dedicated KKT admissibility audit (\texttt{HC-P1-02}, exported to \texttt{data/processed/benchmarks/m1\_9/post\_spinodal\_kkt\_audit.dat}) proves that every serrated point in $h \in [0.400, 0.454]$ is an \textbf{unconverged saddle excursion} (\texttt{conv = False}, $n_\Phi = 60$ limit reached, $\epsilon_{\mathrm{comm}} \in [4.6 \times 10^{-4}, 1.21]$, $\epsilon_{\mathrm{gauge}} \in [3.3 \times 10^{-4}, 3.58 \times 10^{-1}]$). Because the KKT Lagrangian $\mathcal{L}(\Omega, \boldsymbol{\lambda})$ is a min-max saddle (concave in $\boldsymbol{\lambda}$ per Lemma~\ref{lem:m1_3_hf_gradient}), unconverged iterates that violate dual stationarity $g_s^\gamma = 0$ and the kinematic Bloch-ball bound $\|\mathbf{m}_s\|_2 \le 1$ artificially depress the Gaussian energy below the physical ground state. Enforcing strict KKT convergence ($\epsilon_{\mathrm{gauge}} < 10^{-9}$) and Bloch-ball admissibility eliminates all downward spikes: at $h_{c2}^+ = 0.400$, the QSL local minimum annihilates with an unstable saddle, and the trajectory jumps monotonically onto the exact polarized branch $E_{\mathrm{site}}(h) = -h - 0.5$ with $m_{111} = 1.00000$, $W_p = 0.0000$, and $\epsilon_{\mathrm{gauge}} \le 5.6 \times 10^{-15}$.

Furthermore, to test whether an intermediate spatially modulated phase intervenes inside the spinodal hysteresis window $h \in [0.3325, 0.400]$, Milestone M1.9 (\texttt{m1\_9\_t5\_spinodal\_domain\_wall\_coexistence.md}) relaxed a composite spatial domain-wall seed on an $L_x \times L_y = 16 \times 8$ supercell ($N_{\mathrm{sub}} = 128$) at $h = 0.335$. The interface stabilizes into a sharp domain wall of width $\xi_{\mathrm{DW}} \approx 1.5$ lattice spacings with strictly positive surface tension:
\begin{equation}
\label{eq:m1_6_domain_wall_tension}
\sigma_{\mathrm{DW}} \equiv E_{\mathrm{DW}} - \frac{1}{2}\left(E_{\mathrm{QSL}} + E_{\mathrm{POL}}\right) = +0.042\, |K| > 0,
\end{equation}
ruling out spontaneous domain-wall proliferation or intermediate incommensurate stripe phases.

% EDITORIAL QUESTION: Does unit cell enlargement introduce spontaneous nematicity or alter the physical phase diagram?
\section{Supercell Scale Parity and Absence of Spontaneous Symmetry Breaking}
\label{sec:m1_6_supercell_parity}

To verify that the 2-site unit cell ansatz does not artificially constrain the variational search or suppress competing spatial orders, Milestone M1.6 evaluated a 4-site rectangular supercell (\texttt{rectangular\_4site}, 4 sites, 6 bonds, 2 independent hexagonal plaquettes $p_0, p_1$):

\begin{proposition}[Supercell Scale Parity and Stability]
\label{prop:m1_6_supercell_parity}
Across all fields $h \in [0, 2.0]$:
\begin{enumerate}
\item \textbf{Energy Scale Parity:} The ground-state variational energy per physical site matches between 2-site and 4-site representations to within BZ discretization error:
\begin{equation}
\label{eq:m1_6_energy_parity}
|E_{\mathrm{site}}^{(2)}(h) - E_{\mathrm{site}}^{(4)}(h)| \le 2.2 \times 10^{-4}.
\end{equation}
\item \textbf{Flux Sector Uniformity:} The two independent plaquettes maintain identical fluxes: $|W_p(p_0) - W_p(p_1)| = 0.00 \times 10^0$ to machine precision.
\item \textbf{Absence of Spontaneous Nematicity:} Sublattice magnetizations preserve $C_3$ isotropy up to rectangular grid discretization artifacts:
\begin{equation}
\label{eq:m1_6_nematicity_bound}
N_{\mathrm{nem}} \equiv \frac{1}{N_{\mathrm{sub}}} \sum_{s=0}^{N_{\mathrm{sub}}-1} \left( |\sigma_s^x - \sigma_s^y| + |\sigma_s^y - \sigma_s^z| \right) < 5 \times 10^{-4}.
\end{equation}
Here the residual anisotropy of order $4 \times 10^{-4}$ in the 4-site rectangular cell is an artifact of the rectangular $k$-grid breaking the threefold rotational symmetry of the BZ quadrature, rather than spontaneous physical nematic ordering.
\end{enumerate}
\end{proposition}

\begin{figure}[htbp]
\centering
\begin{tikzpicture}
\begin{axis}[
    labbook standard,
    xlabel={Magnetic Field $h/|K|$ ($\mathbf{h} \parallel [111]$)},
    ylabel={Energy per Physical Site $E_{\mathrm{site}}$},
    legend pos=south west,
    xmin=0.0, xmax=2.0,
    ymin=-2.55, ymax=-0.45
]
    \addplot[labbookteal, thick, mark=o, mark size=2pt]
        table[x index=0, y index=1] {../../data/processed/benchmarks/m1_6/supercell_plot.dat};
    \addlegendentry{2-site Hexagonal QSL}

    \addplot[labbookslate, dashed, thick, mark=square*, mark size=1.8pt]
        table[x index=0, y index=2] {../../data/processed/benchmarks/m1_6/supercell_plot.dat};
    \addlegendentry{4-site Rectangular QSL}

    \addplot[labbookcrimson, thick, mark=triangle*, mark size=2pt]
        table[x index=0, y index=3] {../../data/processed/benchmarks/m1_6/supercell_plot.dat};
    \addlegendentry{2-site Polarized}

    \addplot[labbookamber, dashed, thick, mark=diamond*, mark size=1.8pt]
        table[x index=0, y index=4] {../../data/processed/benchmarks/m1_6/supercell_plot.dat};
    \addlegendentry{4-site Polarized}
\end{axis}
\end{tikzpicture}
\caption{Scale parity comparison between the 2-site hexagonal unit cell and the 4-site rectangular cell across $h/|K| \in [0, 2]$. Ground-state energies agree to within $2.2 \times 10^{-4}$ with zero spontaneous nematicity, demonstrating structural robustness of the mean-field landscape.}
\label{fig:m1_6_supercell_parity}
\end{figure}



% EDITORIAL QUESTION: How do non-linear solvers perform under a rigorous statistical ensemble of symmetric perturbations, and what is the exact computational cost breakdown?
\section{Statistical Ensemble Scorecards, Limited-Memory Broyden, and Host Hardware FLOP Profiling}
\label{sec:m1_6_statistical_scorecards}

To establish rigorous statistical guarantees rather than anecdotal single-seed trials, we evaluate an ensemble of $N_{\mathrm{trials}} = 30$ independent $C_3$-symmetric and sublattice-symmetric random seed perturbations per solver across perturbation radii $\delta \in [0.005, 0.500]$ at the intermediate field $h/|K| = 0.20$.

\subsection{Wilson Score Intervals and Statistical Basin Definition}
For each solver and perturbation magnitude $\delta$, let $k$ denote the number of successful convergences out of $N_{\mathrm{trials}}$. The sample point estimate $\hat{p} = k / N_{\mathrm{trials}}$ is accompanied by the Wilson score interval with continuity correction at confidence level $1 - \alpha = 0.95$ ($z = 1.96$):
\begin{equation}
\label{eq:m1_6_wilson_interval}
w^{\pm}(\hat{p}, N) = \frac{\hat{p} + \frac{z^2}{2N} \pm z \sqrt{\frac{\hat{p}(1-\hat{p})}{N} + \frac{z^2}{4N^2}}}{1 + \frac{z^2}{N}}.
\end{equation}
We define the high-reliability basin radius $R_{90}$ as the supremum of perturbation radii $\delta$ for which the lower Wilson bound satisfies $w^- \ge 0.85$ (or point estimate $\hat{p} \ge 0.90$).

\subsection{Limited-Memory Broyden Formulation (L-Broyden)}
In high-dimensional real-space supercells ($D = \dim \mathcal{X} \gg 10^3$), maintaining the full dense approximate inverse Jacobian $H_k \in \mathbb{R}^{D \times D}$ becomes memory-prohibitive ($\mathcal{O}(D^2)$ storage and $\mathcal{O}(D^2)$ FLOPs per step). We formulate the Limited-Memory Broyden (L-Broyden) algorithm with history depth $m$:
\begin{enumerate}
\item Retain the most recent $m$ displacement and residual differences:
\begin{equation}
s_j \equiv X_{j+1} - X_j, \qquad y_j \equiv R_{j+1} - R_j, \qquad j \in \{k-m, \dots, k-1\}.
\end{equation}
\item Implicitly represent the inverse Jacobian update using rank-1 corrections without forming $H_k$:
\begin{equation}
H_{j+1} = H_j + \frac{s_j - H_j y_j}{s_j^T H_j y_j} s_j^T H_j.
\end{equation}
\item By precomputing the forward vectors $u_j \equiv \frac{s_j - H_j y_j}{s_j^T H_j y_j}$ and adjoint vectors $w_j \equiv H_j^T s_j$, the action $H_k R_k$ is evaluated in $\mathcal{O}(m^2 D)$ operations using $\mathcal{O}(mD)$ storage.
\end{enumerate}

\begin{figure}[htbp]
\centering
\begin{tikzpicture}
\begin{axis}[
    labbook standard,
    xlabel={Perturbation Magnitude $\delta$},
    ylabel={Convergence Probability $P_{\mathrm{conv}}(\delta)$},
    legend pos=south west,
    xmin=0.0, xmax=0.52,
    ymin=-0.05, ymax=1.05
]
    \addplot[labbookteal, thick, mark=*, mark size=2pt]
        table[x index=0, y index=2] {../../data/processed/benchmarks/m1_6/basin_survival_plot.dat};
    \addlegendentry{Broyden (Full)}

    \addplot[labbookslate, thick, mark=square*, mark size=1.8pt]
        table[x index=0, y index=1] {../../data/processed/benchmarks/m1_6/basin_survival_plot.dat};
    \addlegendentry{DIIS (Joint)}

    \addplot[labbookpurple, thick, mark=diamond*, mark size=2pt]
        table[x index=0, y index=5] {../../data/processed/benchmarks/m1_6/basin_survival_plot.dat};
    \addlegendentry{JFNK}

    \addplot[labbookamber, thick, mark=triangle*, mark size=2pt]
        table[x index=0, y index=3] {../../data/processed/benchmarks/m1_6/basin_survival_plot.dat};
    \addlegendentry{L-Broyden ($m=10$)}

    \addplot[labbookcrimson, dashed, thick, mark=x, mark size=2.5pt]
        table[x index=0, y index=4] {../../data/processed/benchmarks/m1_6/basin_survival_plot.dat};
    \addlegendentry{Picard ($\alpha=0.5$)}

    % 90% survival threshold line
    \draw[gray, dotted, thick] (axis cs:0,0.9) -- (axis cs:0.52,0.9) node[above, pos=0.25, font=\tiny] {$R_{90}$ threshold};
\end{axis}
\end{tikzpicture}
\caption{Empirical basin survival probability $P_{\mathrm{conv}}(\delta)$ as a function of symmetric random seed perturbation magnitude $\delta$ at $h/|K|=0.20$ ($N_{\mathrm{trials}}=30$ per solver per $\delta$). Full Broyden and JFNK maintain $P_{\mathrm{conv}} \ge 0.967$ across the entire range, whereas Picard rapidly collapses for $\delta \ge 0.05$.}
\label{fig:m1_6_basin_survival}
\end{figure}


\begin{figure}[htbp]
\centering
\begin{tikzpicture}
\begin{axis}[
    labbook standard,
    xmode=log,
    log basis x={2},
    xlabel={Matrix Dimension $M$},
    ylabel={Sustained Throughput [GFLOP/s]},
    legend pos=north west,
    xmin=6, xmax=1500,
    ymin=0, ymax=190
]
    \addplot[labbookteal, thick, mark=*, mark size=2pt]
        table[x index=0, y index=2] {../../data/processed/benchmarks/m1_6/hardware_flops_calibration.dat};
    \addlegendentry{Complex GEMM (\texttt{zgemm})}

    \addplot[labbookslate, thick, mark=square*, mark size=1.8pt]
        table[x index=0, y index=1] {../../data/processed/benchmarks/m1_6/hardware_flops_calibration.dat};
    \addlegendentry{Real GEMM (\texttt{dgemm})}

    \addplot[labbookpurple, thick, mark=diamond*, mark size=2pt]
        table[x index=0, y index=4] {../../data/processed/benchmarks/m1_6/hardware_flops_calibration.dat};
    \addlegendentry{Complex Eigensolver (\texttt{zheevd})}

    \addplot[labbookamber, thick, mark=triangle*, mark size=2pt]
        table[x index=0, y index=3] {../../data/processed/benchmarks/m1_6/hardware_flops_calibration.dat};
    \addlegendentry{Real Symmetric Eigensolver (\texttt{dsyevd})}
\end{axis}
\end{tikzpicture}
\caption{Host CPU hardware throughput calibration on the 12-core Intel i5-13500 platform across matrix dimensions $M \in [8, 1152]$. Peak dense matrix multiplication saturates at $\sim 171\text{ GFLOP/s}$ in double precision, while divide-and-conquer eigensolvers achieve $20\text{--}23\text{ GFLOP/s}$ for large $M$.}
\label{fig:m1_6_flop_calibration}
\end{figure}


\subsection{Empirical Scorecard and Host Hardware Roofline}
The empirical results compiled in Table~\ref{tab:m1_6_solver_scorecard} confirm the fundamental trade-offs:
\begin{enumerate}
\item \textbf{Basin Radius Superiority:} Full Broyden and JFNK achieve unconditional 95\% Wilson survival ($P_{\mathrm{conv}} = 1.000$, Wilson $[0.886, 1.000]$) up to $\delta = 0.35$, yielding $R_{90} \ge 0.50$. In contrast, damped Picard collapses for $\delta \ge 0.05$ ($P_{\mathrm{conv}} = 0.467$).
\item \textbf{Contraction Acceleration:} At $\delta = 0.02$, full Broyden requires median $K_{50} = 27$ iterations compared to $115$ for Picard, representing a $4.26\times$ reduction in evaluation calls.
\item \textbf{L-Broyden Trade-Off:} L-Broyden ($m=10$) achieves $P_{\mathrm{conv}} = 0.967$ at $\delta = 0.005$ and $P_{\mathrm{conv}} = 0.800$ at $\delta = 0.02$, but drops to $0.533$ at $\delta = 0.10$ due to low-rank history truncation in high-curvature regions.
\item \textbf{Hardware FLOP Calibration:} On the Intel i5-13500 CPU, each evaluation $\Phi$ on the $N_k=1024$ mesh consumes exactly $1.355\times 10^6$ FLOPs. Sustained divide-and-conquer eigensolver throughput is $20.2\text{--}23.1\text{ GFLOP/s}$, while dense matrix multiplication reaches $150\text{--}171\text{ GFLOP/s}$. Broyden step overhead is $1.1\times 10^4$ FLOPs ($<0.8\%$ of $\Phi$), confirming that eigensolvers dominate total execution time.
\end{enumerate}

\begin{table}[htbp]
\centering
\small
\begin{tabular}{lcccccc}
\toprule
\textbf{Solver} & \textbf{$R_{90}$ Radius} & \textbf{$R_{50}$ Radius} & \textbf{$K_{50}$ ($\delta=0.02$)} & \textbf{Time ($\delta=0.02$)} & \textbf{Storage} & \textbf{Step FLOPs} \\
\midrule
Full Broyden   & $\ge 0.500$ & $\ge 0.500$ & 27.0 & $327\text{ ms}$ & $\mathcal{O}(D^2)$ & $2 D^2$ \\
DIIS (Joint)   & $0.350$     & $0.500$     & 34.0 & $242\text{ ms}$ & $\mathcal{O}(mD)$   & $\mathcal{O}(m^2 D + m^3)$ \\
JFNK           & $\ge 0.500$ & $\ge 0.500$ & 39.0 & $400\text{ ms}$ & $\mathcal{O}(k_{\mathrm{kry}} D)$ & $\mathcal{O}(k_{\mathrm{kry}} \Phi)$ \\
L-Broyden ($m=10$) & $0.010$ & $0.100$     & 55.0 & $465\text{ ms}$ & $\mathcal{O}(mD)$   & $\mathcal{O}(m^2 D)$ \\
Picard ($\alpha=0.5$) & $0.010$ & $0.020$ & 115.0 & $1133\text{ ms}$ & $\mathcal{O}(D)$   & $\mathcal{O}(D)$ \\
\bottomrule
\end{tabular}
\caption{Comprehensive statistical solver scorecard at $h/|K|=0.20$ across $N_{\mathrm{trials}}=30$ symmetric random seed trials. Full Broyden and DIIS form the optimal Pareto efficiency frontier for Bloch supercells.}
\label{tab:m1_6_solver_scorecard}
\end{table}
